% Part: first-order-logic
% Chapter: sequent-calculus
% Section: provability-consistency

\documentclass[../../../include/open-logic-section]{subfiles}

\begin{document}

\iftag{FOL}
      {\olfileid[fr]{fol}{seq}{prv}}
      {\olfileid[fr]{pl}{seq}{prv}}

\olsection{\usetoken{S}{derivability} et cohérence}

Établissons plusieurs propriétés de la relation de !!{derivability}.
Outre leur intérêt propre, elles interviendront toutes dans la
démonstration du théorème de complétude.

\begin{prop}\ollabel{prop:provability-contr}
Si $\Gamma \Proves !A$ et si $\Gamma \cup \{!A\}$ est incohérent,
alors $\Gamma$ est incohérent.
\end{prop}

\begin{proof}
Il existe des parties finies $\Gamma_0 \subseteq \Gamma$ et
$\Gamma_1 \subseteq \Gamma$ telles que $\Log{LK}$ !!{derive}s
$\Gamma_0 \Sequent !A$ et $!A, \Gamma_1 \Sequent \quad$,
en ajoutant au besoin $!A$ par affaiblissement dans la seconde.
Notons $\pi_0$ et $\pi_1$ les !!{derivation}s respectives de ces
deux séquents dans $\Log{LK}$.\footnote{Note de l'édition française :
la présentation de $\pi_1$ est uniformisée avec l'arbre qui suit ;
la source intervertit $!A$ et $\Gamma_1$ à cet endroit.}
Nous pouvons alors !!{derive} :
\begin{prooftree}
\AxiomC{}
\RightLabel{$\pi_0$}
\Deduce$ \Gamma_0 \fCenter !A $
\AxiomC{}
\RightLabel{$\pi_1$}
\Deduce$!A, \Gamma_1 \fCenter $
\RightLabel{\Cut}
\BinaryInf$ \Gamma_0,\Gamma_1 \fCenter $
\end{prooftree}

Comme $\Gamma_0 \subseteq \Gamma$ et $\Gamma_1 \subseteq \Gamma$,
on a $\Gamma_0 \cup \Gamma_1 \subseteq \Gamma$ ; ainsi,
$\Gamma$ est incohérent.
\end{proof}

\begin{prop}
\ollabel{prop:prov-incons}
$\Gamma \Proves !A$ si et seulement si
$\Gamma \cup \{\lnot !A\}$ est incohérent.
\end{prop}

\begin{proof}
Supposons d'abord $\Gamma \Proves !A$. Il existe une partie finie
$\Gamma_0 \subseteq \Gamma$ et une !!{derivation}~$\pi_0$ de
$\Gamma_0 \Sequent !A$. En ajoutant une inférence $\LeftR{\lnot}$,
on obtient une !!{derivation} de
$\lnot !A, \Gamma_0 \Sequent \quad$. L'ensemble
$\Gamma \cup \{\lnot !A\}$ est donc incohérent.\footnote{Note de
l'édition française : dans les deux sens de cette démonstration,
la source emploie directement $\Gamma$ dans les séquents.
Les témoins finis $\Gamma_0 \subseteq \Gamma$ sont explicités ici,
car l'ensemble $\Gamma$ peut être infini.}

Réciproquement, si $\Gamma \cup \{\lnot !A\}$ est incohérent,
il existe une partie finie $\Gamma_0 \subseteq \Gamma$ et une
!!{derivation}~$\pi_1$ de $\lnot !A, \Gamma_0 \Sequent \quad$,
en ajoutant au besoin $\lnot !A$ par affaiblissement.
L'arbre suivant est une !!{derivation} de $\Gamma_0 \Sequent !A$ :
\begin{prooftree}
  \Axiom$!A \fCenter !A$
  \RightLabel{\RightR{\lnot}}
  \UnaryInf$\fCenter !A, \lnot !A$
  \AxiomC{}
  \RightLabel{$\pi_1$}
  \Deduce$\lnot !A, \Gamma_0 \fCenter$
  \RightLabel{\Cut}
  \BinaryInf$\Gamma_0 \fCenter !A$
\end{prooftree}
Puisque $\Gamma_0 \subseteq \Gamma$, il en résulte
$\Gamma \Proves !A$.
\end{proof}

\begin{prob}
Démontrer que $\Gamma \Proves \lnot !A$ si et seulement si
$\Gamma \cup \{!A\}$ est incohérent.
\end{prob}

\begin{prop}\ollabel{prop:explicit-inc}
Si $\Gamma \Proves !A$ et $\lnot !A \in \Gamma$, alors
$\Gamma$ est incohérent.
\end{prop}

\begin{proof}
Supposons $\Gamma \Proves !A$ et $\lnot !A \in \Gamma$.
Il existe une partie finie $\Gamma_0 \subseteq \Gamma$ et une
!!{derivation}~$\pi$ du séquent $\Gamma_0 \Sequent !A$.
Le séquent $\lnot !A, \Gamma_0 \Sequent \quad$ est lui aussi
!!{derivable} : l'arbre suivant donne son antécédent dans l'ordre
inverse, que l'on peut rétablir par des échanges.
  \begin{prooftree}
    \AxiomC{}
    \RightLabel{$\pi$}
    \Deduce$\Gamma_0 \fCenter !A$
    \Axiom$!A \fCenter !A$
    \RightLabel{\LeftR{\lnot}}
    \UnaryInf$\lnot !A, !A \fCenter$
    \RightLabel{\LeftR{\Exchange}}
    \UnaryInf$!A, \lnot !A \fCenter$
    \RightLabel{\Cut}
    \BinaryInf$\Gamma_0, \lnot !A \fCenter$
  \end{prooftree}
Comme $\lnot !A \in \Gamma$ et $\Gamma_0 \subseteq \Gamma$,
cela montre que $\Gamma$ est incohérent.
\end{proof}

\begin{prop}\ollabel{prop:provability-exhaustive}
Si $\Gamma \cup \{!A\}$ et $\Gamma \cup \{\lnot !A\}$ sont
tous deux incohérents, alors $\Gamma$ est incohérent.
\end{prop}

\begin{proof}
Il existe des parties finies $\Gamma_0 \subseteq \Gamma$ et
$\Gamma_1 \subseteq \Gamma$, ainsi que des !!{derivation}s
$\pi_0$ et $\pi_1$ dans $\Log{LK}$, respectivement de
$!A, \Gamma_0 \Sequent \quad$ et de
$\lnot !A, \Gamma_1 \Sequent \quad$.
On ajoute au besoin l'énoncé distingué par affaiblissement.
Nous pouvons alors !!{derive} :
\begin{prooftree}
\AxiomC{}
\RightLabel{$\pi_0$}
\Deduce$ !A, \Gamma_0 \fCenter $
\RightLabel{\RightR{\lnot}}
\UnaryInf$ \Gamma_0 \fCenter \lnot !A$
\AxiomC{}
\RightLabel{$\pi_1$}
\Deduce$\lnot !A, \Gamma_1 \fCenter  $
\RightLabel{\Cut}
\BinaryInf$ \Gamma_0, \Gamma_1 \fCenter $
\end{prooftree}
Comme $\Gamma_0 \subseteq \Gamma$ et $\Gamma_1 \subseteq \Gamma$,
on a $\Gamma_0 \cup \Gamma_1 \subseteq \Gamma$.
Par conséquent, $\Gamma$ est incohérent.
\end{proof}

\end{document}
